% RESEARCH_DRAFT_SHA256: 579e2a515090ad1e9fa88e1f16b4cf446fb3236b2ff1d9cb975400ade6620660
% Layout and source-reference wrappers expanded for the web renderer.
\documentclass[11pt]{article}
\usepackage[margin=25mm]{geometry}
\usepackage{fontspec,unicode-math}
\setmainfont{Latin Modern Roman}
\setmathfont{Latin Modern Math}
\usepackage{amsmath,amssymb,mathtools,xurl,hyperref}
\hypersetup{colorlinks=true,urlcolor=blue,linkcolor=blue}
\setlength{\parindent}{0pt}
\setlength{\parskip}{5pt}
\setlength{\emergencystretch}{3em}
\newcommand{\sref}[2]{\hyperlink{source:#1:#2}{[S#2]}}
\newcommand{\cataloguescope}[1]{\paragraph{Source question.}#1}
\newcommand{\floor}[1]{\left\lfloor#1\right\rfloor}
\begin{document}

\section*{Fixed cubic root finding: exact easy cases and conditional ternary correction}
\noindent catalogue entry; catalogue entry. Research attempt dated 11 October 2026.
\paragraph{Result in scope.}This manuscript contains a source-specific attempt and complete proofs of the stated auxiliary calculations. It makes no new solution or novelty claim for the entire target.
\section{Original problem and preserved definition}
The following definition and source list are reproduced from the repository. Any earlier inspection dates and reported literature claims in this preserved text describe that source record, not additional verification by this attempt.
% BEGIN PRESERVED DEFINITION
\subsection{Definitions and mathematical statement}

\paragraph{Shared notation.}Standard number systems are denoted by $\mathbb N,\mathbb Z,\mathbb Q,\mathbb R,\mathbb C$; index ranges, computational models and specialized notation follow the cited problem source and the definitions below.


Polynomial time means a bound polynomial in the finite input length; for an integer input, its binary length is proportional to its logarithm. A root modulo p of a fixed polynomial f is an element $a\in\mathbb F_p$ with $f(a)=0$. The sample polynomial is $x^3-2$. Since f is fixed, polynomial time is measured in the binary length of the prime p; only roots that exist can be output.
\paragraph{Statement recorded in the supplied source.}
Question 18. Can you compute roots modulo p of a fixed polynomial (` a la Schoof-Pila) in polynomial time, like x3 − 2? (Schoof) 30 \href{https://aimath.org/WWN/primesinp/primesinp.pdf}{[S2]}
\paragraph{Scope qualification.}
The source asks in the context of deterministic Schoof--Pila-style methods; it does not claim every prime yields a root.
\paragraph{Source status.}
The supplied catalogue labels this record partially solved. Its source-status metadata was checked against the supplied snapshot on 2026-10-11. This import does not independently establish current open status or mathematical correctness.
\subsection{Short English statement}
Can roots of a fixed polynomial, such as x cubed minus two, be found modulo primes in polynomial time?
\subsection{Sources}
\begin{itemize}
\item[S1] Ulam AI and the UnsolvedMath Contributors, UnsolvedMath: A Curated Collection of Open Mathematics Problems, catalogue entry. Catalogue curation and metadata: CC BY 4.0. The supplied source question is retained; definitions and the English summary are editorial additions.
\url{https://huggingface.co/datasets/ulamai/UnsolvedMath}.
\item[S2] AIM Workshop Problem Lists, original problem source linked by the catalogue
\url{https://aimath.org/WWN/primesinp/primesinp.pdf}.
\end{itemize}
% END PRESERVED DEFINITION
\subsection{Verbatim supplied catalogue question}
Question 18. Can you compute roots modulo p of a fixed polynomial (` a la Schoof-Pila) in polynomial time, like x3 − 2? (Schoof) 30

\section{Target and attempted exponent route}
Question 18 asks for a deterministic polynomial-time algorithm, for a
fixed polynomial such as $X^3-2$, finding its roots modulo a varying
prime $p$. Here polynomial time means polynomial in $\log p$ with
the polynomial fixed. We prove unconditional elementary cases and a
complete remaining root-extraction algorithm conditional on a supplied
cubic nonresidue. We do not supply that nonresidue deterministically
in the stated time for every prime, and do not infer such a construction
from the source's complex-multiplication remark.

\paragraph{The finite-field group fact used below.}
The multiplicative group of a finite field is cyclic. A short proof
is useful here. Let $e$ be the least common multiple of the orders
of its elements. For each prime power dividing $e$ to its exact
exponent, some element order contains that prime power; taking an
appropriate power of the element gives precisely that order. Products
of these elements have order the product of their pairwise coprime
orders: if such a product raised to $k$ is one, isolating a factor
and raising to the product of the other orders shows its order divides
$k$ by coprimality. Thus an element of order $e$ exists. All group
elements are roots of $X^e-1$, so their number is at most $e$;
the element's distinct powers give the opposite inequality. It
therefore generates the entire group. The root bound used here
follows by repeatedly dividing a polynomial by $X-a$ for a root.

\paragraph{Unconditional branches.}
For $p=2$ the only root is zero; for $p=3$ it is two, by direct
evaluation of the field elements. If $p\equiv2\pmod3$, set
$e=(2p-1)/3$. Then $3e=2(p-1)+1$, so
$(2^e)^3=2$ by the order-$p-1$ multiplicative group identity.
The cube map is bijective because multiplication by three permutes
exponents modulo $p-1$. Thus this gives all roots, namely the unique
one, using binary exponentiation in polynomial time.

For $p\equiv1\pmod3$, writing a generator as $w$ shows that
$a\ne0$ is a cube exactly when $a^{(p-1)/3}=1$: both conditions
say that its exponent relative to $w$ is divisible by three.
This tests whether $2$ has any roots with polynomially many bit
operations. If $p-1=3m$ with $3\nmid m$ and $2$ passes the
test, choose $e$ inverse to three modulo $m$ (if $m=1$ set $e=0$).
Then $2^m=1$ and $(2^e)^3=2$, so one root is found.
Finding all three additionally needs a nontrivial cube root of unity;
the one-root calculation alone has not constructed it.

\paragraph{A complete conditional algorithm for the other cases.}
Suppose $p-1=3^s m$, $s\ge1$, $3\nmid m$, and a noncube
$z\in\mathbb F_p^*$ is supplied. Set $c=z^m$. Relative to
a generator, the exponent of $z$ is not divisible by three,
so $c$ has exact order $3^s$. Assume $a=2$ passes the cube test.
Choose $e$ with $3e\equiv1\pmod m$, using $e=0$ when $m=1$,
and put $b=a^e$, $u=a/b^3$. The exponent $1-3e$ is a multiple
of $m$, so $u$ belongs to the cyclic subgroup generated by $c$.
Moreover $a$ and $b^3$ are cubes, so $u$ is a cube. In this
subgroup that says $u=c^{3j}$ for a unique $j$ modulo $3^{s-1}$:
if $c=w^{mk}$ with $3\nmid k$, then $c^j$ is a cube
precisely when $3\mid j$.

This discrete logarithm in a power-of-three group can be computed
digit by digit. Put $h=c^3$, of order $3^{s-1}$. If $s=1$,
then $u=1$ and set $j=0$. Otherwise start $J=0$ and, for
$k=0,\ldots,s-2$, form
\[
R=u h^{-J},\qquad D=R^{3^{s-2-k}},\qquad
\omega=h^{3^{s-2}}.
\]
Choose the unique $d_k\in\{0,1,2\}$ for which $D=\omega^{d_k}$,
and replace $J$ by $J+d_k3^k$. To verify this procedure, write
the true $j$ in base three. If its first $k$ digits agree with
$J$, then $j-J=d_k3^k+3^{k+1}v$. Consequently
$R^{3^{s-2-k}}=h^{d_k3^{s-2}}$, since the remaining exponent
is divisible by $3^{s-1}$. The three powers of $\omega$ are
distinct because its order is three. This proves the inductive
digit step, and at the end $u=h^J$. Therefore
\[
x=b c^J\quad\hbox{satisfies}\quad x^3=b^3u=a.
\]
All roots are $x,x\xi,x\xi^2$, where
$\xi=c^{3^{s-1}}$ has order three. They are distinct and exhaust
the roots because a cubic has at most three roots. There are
$s\le\log_3(p-1)$ digit steps, each using modular powers with
$O(\log p)$-bit exponents and arithmetic, so the conditional
algorithm is polynomial in $\log p$.

\section{First missing unconditional step and audit}
The primary Question 18 was inspected. The proof locates the precise
obstruction in this route: an unconditional deterministic polynomial-time
construction of a cubic nonresidue, or another mechanism avoiding it.
Trying $z=1,2,3,\ldots$ and testing each eventually succeeds, but no
polynomial-in-$\log p$ bound for this search is proved here.
Existence follows from cyclicity, since two thirds of nonzero elements
are noncubes; that density does not bound their first position in a
deterministic ordered search. Random sampling is not the requested
deterministic algorithm. The audit distinguishes finding one root
from all roots and handles $m=1$, $s=1$ and the primes two and
three separately. The complete argument is algebraic; finite exact
checks described below test implementation but not the unresolved bound.
Exact Python integer arithmetic checked every supplied noncube for each
prime below 500 congruent to one modulo three for which 2 is a cube.
There were 2,664 input pairs across 15 primes. In every case all three
outputs agreed with direct enumeration of $x^3=2$ in the field.
The recorded prime list and case count are in the accompanying audit.

\section{Process and contributions}
Mathematics and model-assisted adversarial self-audit: OpenAI gpt-6-astra, ultra reasoning, as dispatched for this research batch; exact serving revision not exposed. Codex (GPT-6; exact serving revision not exposed) materially assisted with orchestration and format review. No human mathematical authorship, mathematical edits, or independent human proof checking is asserted.
Contribution prompt A was applied to the full definition and supplied source row, then prompt B was applied adversarially to the derived results. The item-specific audit above records the scope checks and rejected extrapolations. Local assembly and structural checks do not establish mathematical correctness. The source catalogue status is preserved in TOP\_PROBLEM and is separate from this attempt status.
\paragraph{Integration review.} Codex (GPT-6; exact serving revision not exposed) reviewed the written arguments and expanded layout and source-reference wrappers. Repository metadata and the production web renderer were checked using Python 3.13 and JavaScript with MathJax 3.2.2. This model-assisted review is not independent human proof verification; no PDF compilation or formal proof checking is asserted.

\section{Completion Estimate}
\noindent\textbf{COMPLETION: 0\%}
The proved material is a rederivation of known or auxiliary facts rather than a substantiated new advance on the original target. Unconditional congruence branches and a complete conditional ternary root algorithm are proved; a deterministic polynomial-time nonresidue construction or replacement remains missing.
\section{References}
The complete original Sources subsection above is retained, including catalogue attribution and licence notices. The additional list records primary sources actually consulted for this attempt on 11 October 2026; uncited source links reproduced above were not all rechecked.
\begin{itemize}
\item American Institute of Mathematics, Future directions in algorithmic number theory, workshop 24--28 March 2003, report version 30 April 2003; original question, relevant lecture context and adjoining remarks inspected on 11 October 2026. \url{https://aimath.org/WWN/primesinp/primesinp.pdf}.
\end{itemize}
\section{Final status}
\textbf{Unresolved in this attempt.} Unconditional congruence branches and a complete conditional ternary root algorithm are proved; a deterministic polynomial-time nonresidue construction or replacement remains missing.
\end{document}
